% !PRINTED 570
\weberpage{570}
%<<CONT
actually to decompose any arbitrarily given natural prime $p$ into its
prime factors, and thereby actually to represent all prime ideals of
the field $\Omega$, or at least representatives of all the
prime ideals\footnote{We follow here a paper by \gesperrt{Hensel}: ``Untersuchung der
Fundamentalgleichung einer Gattung für eine reelle Primzahl als Modul
und Bestimmung der Theiler der Discriminante'' (\gesperrt{Crelle}'s Journ., Bd. 113).
Also to be mentioned are the other papers by \gesperrt{Hensel}, Ibid.,
Bd. 101, 104, 105, 113.}.

Let $\mathfrak p$ be an arbitrary prime ideal of degree $f$, so that
\weq{1}{N(\mathfrak p) = p^f}
where $p$ denotes a natural prime which is divisible by the
prime factor $\mathfrak p$. $f$ is a positive exponent
$\leqq n$. The least integral rational number that is divisible by
$\mathfrak p$ is $p$, and every other integral rational number that is
divisible by $\mathfrak p$ is therefore also divisible by $p$
(\S.\,138).

We must now also consider congruences with the modulus $\mathfrak p$
between integral functions of arbitrary variables with coefficients in
$\mathfrak o$, and we remark that by \S.\,143, 2. two such functions
are \gesperrt{congruent then and only then, when the corresponding
coefficients are congruent.}

We denote the field of rational numbers by $R$, and accordingly we
also call a function with rational coefficients a \gesperrt{function
in} $R$.

The basis form of $\mathfrak o$
\weq{2}{\tau = \omega_1 t_1 + \omega_2 t_2 + \cdots + \omega_n t_n}
satisfies, as we have seen in \S.\,146, an equation of degree
$n$, $F(\tau) = 0$, whose coefficients are integral
rational functions of $t_1, t_2, \ldots, t_n$. Thus $F(\tau)$ is in
any case divisible by $\mathfrak p$, and from this it follows that
there are integral functions $\Phi(t)$ in $R$, which besides $t$
may contain any other variables, and which by the substitution
$t = \tau$ pass over into functionals in $\mathfrak o$ divisible by
$\mathfrak p$, and which therefore satisfy the congruence
\weq{3}{\Phi(\tau) \equiv 0\ (\mathrm{mod}\ \mathfrak p)}
and so we shall be able to say that $\tau$ \gesperrt{is a root of the
congruence}
\weq{4}{\Phi(t) \equiv 0\ (\mathrm{mod}\ \mathfrak p).}

The function $\Phi$ will certainly have to contain the variables
$t_1, t_2, \ldots, t_n$; but it can also contain still other
%CONT>>
